% ═══════════════════════════════════════════════════════════════
% MT-03a-ML: Minitest Las habitaciones - MUSTERLÖSUNG
% Spanisch Klasse 7 | Sequenz 3: Mi casa
% ═══════════════════════════════════════════════════════════════

\documentclass[11pt, a4paper]{article}

\usepackage[utf8]{inputenc}
\usepackage[T1]{fontenc}
\usepackage[ngerman]{babel}
\usepackage{helvet}
\renewcommand{\familydefault}{\sfdefault}

\usepackage[margin=2cm]{geometry}
\usepackage{enumitem}
\usepackage{tabularx}
\usepackage{booktabs}

\usepackage{xcolor}
\usepackage{tcolorbox}
\tcbuselibrary{skins}

\usepackage{fancyhdr}
\usepackage{lastpage}
\usepackage{needspace}

% FARBEN
\definecolor{spanisch-color}{HTML}{c41e3a}
\definecolor{grau-color}{HTML}{888888}
\definecolor{hellgrau-color}{HTML}{f5f5f5}
\definecolor{loesung-color}{HTML}{c62828}

\colorlet{spanisch}{spanisch-color}
\colorlet{grau}{grau-color}
\colorlet{hellgrau}{hellgrau-color}
\colorlet{loesung}{loesung-color}

\newcommand{\fachfarbe}{spanisch}

\newcommand{\thema}{Minitest: Las habitaciones}
\newcommand{\fach}{Spanisch}
\newcommand{\klasse}{7}
\newcommand{\maxpunkte}{26}
\newcommand{\zeit}{15 Minuten}
\newcommand{\datum}{\today}

\pagestyle{fancy}
\fancyhf{}
\renewcommand{\headrulewidth}{0.4pt}
\renewcommand{\footrulewidth}{0.4pt}

\lhead{\fach{} -- Klasse \klasse}
\chead{\textcolor{loesung}{\textbf{MT-03a MUSTERLÖSUNG}}}
\rhead{\datum}

\lfoot{\zeit}
\cfoot{}
\rfoot{Seite \thepage\ von \pageref{LastPage}}

\newcommand{\punktebox}[1]{\hfill\fbox{\small \_/#1}}
\newcommand{\luecke}[1]{\rule{#1}{0.4pt}}
\newcommand{\lsg}[1]{\textcolor{loesung}{\textbf{#1}}}

\newtcolorbox{infobox}{
    colback=hellgrau,
    colframe=\fachfarbe,
    boxrule=1pt,
    arc=0pt,
    left=8pt, right=8pt, top=8pt, bottom=8pt
}

\newenvironment{aufgabe}[1]{%
    \needspace{5\baselineskip}%
    \nopagebreak[4]%
    \vspace{10pt}
    \noindent\textbf{\textcolor{\fachfarbe}{Aufgabe #1}}
    \nopagebreak[4]%
    \vspace{5pt}
    \nopagebreak[4]%
    \begin{list}{}{\leftmargin=0pt}
    \item[]
}{%
    \end{list}
}

\newcommand{\es}[1]{\textcolor{\fachfarbe}{\textbf{#1}}}

\begin{document}

% ═══════════════════════════════════════════════════════════════
% KOPFBEREICH
% ═══════════════════════════════════════════════════════════════

\begin{center}
    {\Large\bfseries\textcolor{\fachfarbe}{\thema}}\\[0.3em]
    {\large\textcolor{loesung}{\textbf{--- MUSTERLÖSUNG ---}}}
\end{center}

\vspace{5pt}

\noindent
\begin{tabularx}{\textwidth}{|X|X|X|}
\hline
\textbf{Name:} \lsg{Musterschüler/in} &
\textbf{Klasse:} \klasse &
\textbf{Datum:} \datum \\
\hline
\end{tabularx}

\vspace{10pt}

\begin{infobox}
\textbf{Zeit:} \zeit{} | \textbf{Punkte:} \maxpunkte{} BE | \textbf{Hilfsmittel:} keine
\end{infobox}

\vspace{15pt}

% ═══════════════════════════════════════════════════════════════
% AUFGABE 1: Zuordnung Räume (8 BE)
% ═══════════════════════════════════════════════════════════════

\begin{aufgabe}{1 -- Räume zuordnen}
Verbinde die spanischen Räume mit der deutschen Übersetzung. \punktebox{8}
\end{aufgabe}

\begin{center}
\begin{tabular}{rclp{4cm}}
\es{la cocina} & \lsg{---} & das Badezimmer & \lsg{$\rightarrow$ die Küche} \\[0.8em]
\es{el salón} & \lsg{---} & die Küche & \lsg{$\rightarrow$ das Wohnzimmer} \\[0.8em]
\es{el dormitorio} & \lsg{---} & der Garten & \lsg{$\rightarrow$ das Schlafzimmer} \\[0.8em]
\es{el baño} & \lsg{---} & das Wohnzimmer & \lsg{$\rightarrow$ das Badezimmer} \\[0.8em]
\es{el jardín} & \lsg{---} & die Garage & \lsg{$\rightarrow$ der Garten} \\[0.8em]
\es{el garaje} & \lsg{---} & das Schlafzimmer & \lsg{$\rightarrow$ die Garage} \\[0.8em]
\es{la terraza} & \lsg{---} & das Esszimmer & \lsg{$\rightarrow$ die Terrasse} \\[0.8em]
\es{el comedor} & \lsg{---} & die Terrasse & \lsg{$\rightarrow$ das Esszimmer} \\
\end{tabular}
\end{center}

\vspace{15pt}

% ═══════════════════════════════════════════════════════════════
% AUFGABE 2: Artikel einsetzen (8 BE)
% ═══════════════════════════════════════════════════════════════

\begin{aufgabe}{2 -- Artikel einsetzen}
Setze den richtigen Artikel ein: \es{el} oder \es{la}? \punktebox{8}
\end{aufgabe}

\noindent
a) \lsg{la} cocina \hfill
b) \lsg{el} salón \\[0.8em]
c) \lsg{el} dormitorio \hfill
d) \lsg{el} baño \\[0.8em]
e) \lsg{la} terraza \hfill
f) \lsg{el} jardín \\[0.8em]
g) \lsg{el} garaje \hfill
h) \lsg{el} comedor \\

\vspace{15pt}

% ═══════════════════════════════════════════════════════════════
% AUFGABE 3: Sätze vervollständigen (10 BE)
% ═══════════════════════════════════════════════════════════════

\begin{aufgabe}{3 -- Sätze vervollständigen}
Ergänze die fehlenden Wörter. Achte auf den Artikel! \punktebox{10}
\end{aufgabe}

\noindent
a) Yo duermo en \lsg{el dormitorio}. \hfill (2P) \\
\hspace*{0.5cm}\textit{(Ich schlafe im Schlafzimmer.)}

\vspace{0.8em}

\noindent
b) Mi madre cocina en \lsg{la cocina}. \hfill (2P) \\
\hspace*{0.5cm}\textit{(Meine Mutter kocht in der Küche.)}

\vspace{0.8em}

\noindent
c) El coche está en \lsg{el garaje}. \hfill (2P) \\
\hspace*{0.5cm}\textit{(Das Auto ist in der Garage.)}

\vspace{0.8em}

\noindent
d) Vemos la tele en \lsg{el salón}. \hfill (2P) \\
\hspace*{0.5cm}\textit{(Wir schauen Fernsehen im Wohnzimmer.)}

\vspace{0.8em}

\noindent
e) Las flores están en \lsg{el jardín}. \hfill (2P) \\
\hspace*{0.5cm}\textit{(Die Blumen sind im Garten.)}

% ═══════════════════════════════════════════════════════════════
% PUNKTEÜBERSICHT
% ═══════════════════════════════════════════════════════════════

\vspace{25pt}
\hrule
\vspace{10pt}

\noindent
\begin{tabularx}{\textwidth}{|X|c|c|c||c|}
\hline
\textbf{Aufgabe} & 1 & 2 & 3 & \textbf{Gesamt} \\
\hline
\textbf{max. Punkte} & 8 & 8 & 10 & \textbf{26} \\
\hline
\textbf{erreicht} & \lsg{8} & \lsg{8} & \lsg{10} & \lsg{26} \\
\hline
\end{tabularx}

\vspace{10pt}

\begin{flushright}
\textbf{Note:} \lsg{1 (14 NP)}
\end{flushright}

\vspace{1em}

\textbf{14-NP-Skala (Sek I):}

\begin{tabular}{|c|c|c|c|c|c|c|c|c|c|c|c|c|c|c|}
\hline
\textbf{NP} & 14 & 13 & 12 & 11 & 10 & 9 & 8 & 7 & 6 & 5 & 4 & 3 & 2 & 1 \\
\hline
\textbf{ab BE} & 26 & 25 & 24 & 22 & 21 & 19 & 17 & 16 & 14 & 12 & 10 & 9 & 7 & 5 \\
\hline
\end{tabular}

\end{document}
